When an international oil company divests an aging producing asset to an indigenous operator in Nigeria, the wells, facilities, and data transfer at completion, but the operating organization that held the asset within safe and productive bounds does not. Outcomes diverged sharply: some acquirers sustained production and integrity, others suffered failures, including an uncontrolled blowout. That divergence is most consistently explained by capability-related factors.

The evidence is six divestment cases in a discriminating-case design. The critical experiment is the OML 29 transfer to Aiteo, a well-funded operator that raised production and then could not contain a major well-control event; the remaining transfers provide cross-case validation. Each case is read on two outcome dimensions, production performance and operational integrity, against five competing explanations.

Capital constraints are rejected as a sufficient cause, fiscal regime effects on grounds of timing; regulatory inadequacy is an enabling condition, and political economy explains access and commercial survival but not operational integrity. Capability-related factors are the explanation most consistent with the combined pattern. A within-asset comparison on OML 30, where operatorship changed while the asset was held fixed and performance improved, strengthens that reading against an asset-quality account.

Such a divestment is a sociotechnical control failure that no voluntary transfer corrects. No party holds both the knowledge of what is lost and the incentive to transfer it, and much of the cost of failure falls on host communities and the state. The contribution is a Capability Assurance Framework, proposed as a mandatory condition on high-risk producing-asset transfers in which seller disclosure, buyer absorption, and regulator verification are mutually constraining, operationalized as six recommendations that re-specify instruments Nigeria already holds. The framework is warranted by the documented conditions rather than proven by implementation. It makes the operating capability a transfer requires observable, testable, and enforceable at the point of transfer.
